\section{Auswertung des Jahresabschlusses 2025}
\label{sec:abschluss}

Auch dieser Teil ist beschreibend. Jede Kennzahl nennt zuerst ihre Gleichung
und dann den Rechenweg; ein Ergebnis ohne seine Formel ist nicht pr\"ufbar.

\subsection{Umsatz}

Der Umsatz stieg von \MioUSDexakt{\UmsatzVJ} auf \MioUSDexakt{\Umsatz} und
damit um \Pct{\UmsatzDelta}.\quelleFS{2025}{7}\quelleMerken{s6fsxcacfxh} Er ist ein Rekordwert.

\subsection{Produktion und Marge je Unze}
\label{sec:margejeunze}

Der Umsatzzuwachs stammt nicht aus der Menge. Die Goldproduktion ging von
\Unzen{\ProduktionXXIV} auf \Unzen{\ProduktionXXV} zur\"uck, also um
\PctBetrag{\ProduktionDelta}; der erzielte Goldpreis stieg zugleich von
\JeUnze{\GoldpreisVJ} auf \JeUnze{\Goldpreis} oder um
\Pct{\GoldpreisDelta}.\quelleMDA{2025}{4}\quelleMerken{s7mdaxcacfxe}

Der eigentliche Fr\"uhindikator eines Goldproduzenten ist deshalb weder der
Umsatz noch die Menge, sondern die Marge je Unze: der erzielte Preis abz\"uglich
der All-in Sustaining Costs.

\begin{equation}
\label{eq:marge-je-unze}
\text{Marge je Unze} = \text{erzielter Goldpreis} - \text{AISC}
\end{equation}

Sie stieg von \JeUnze{\MargeJeUnzeVJ} auf \JeUnze{\MargeJeUnze}, also um
\Pct{\MargeJeUnzeDelta} -- bei sinkender Menge und steigenden St\"uckkosten.
Die Stellgr\"o{\ss}e liegt damit ausserhalb des Unternehmens.

Die St\"uckkosten stiegen von \JeUnze{\AISCVJneu} auf \JeUnze{\AISCXXV}
oder um \Pct{\AISCDelta}.\quelleErneut{s7mdaxcacfxe}

\subsubsection*{Eine Neuberechnung, die benannt geh\"ort}

Der Vorjahreswert der All-in Sustaining Costs lautet im Bericht 2025
\JeUnze{\AISCVJneu}, im Bericht 2024 dagegen
\JeUnze{\AISCXXIV}.\quelleMDA{2024}{4}\quelleMerken{s7mdaxcacexe} Die Differenz von
\JeUnze{\AISCNeuberechnung} ist kein Druckfehler: Das Unternehmen hat die
Vergleichszahlen bereinigt und die Marktbewertungseffekte der
aktienbasierten Verg\"utung aus der Kennzahl herausgenommen.\quelleMDA{2025}{41}\quelleMerken{s7mdaxcacfxeb}

\bk{Die Neuberechnung ist erkl\"art und damit unproblematisch. Sie hat aber
eine Folge, die der Bericht nicht zieht: Wer eine Kostenreihe \"uber mehrere
Jahresberichte hinweg zusammenstellt, mischt ohne es zu merken zwei
Rechenweisen. Tabelle~\ref{tab:reihe} f\"uhrt jede Zahl deshalb so, wie der
Bericht ihres eigenen Jahres sie ausgewiesen hat.}

\subsection{Betriebsergebnis und operatives EBITDA}

Das Betriebsergebnis betr\"agt \MioUSDexakt{\BetriebsErgebnis} nach
\MioUSDexakt{\BetriebsErgebnisVJ}.\quelleFS{2025}{7}\quelleMerken{s7fsxcacfxh} Es enth\"alt allerdings
eine Wertaufholung auf das Bergbauverm\"ogen von
\MioUSDexakt{\Wertaufholung}, die kein Zahlungsstrom ist.

Das Dokument f\"uhrt deshalb als operative Gr\"o{\ss}e das EBITDA ohne diese
Wertaufholung:

\begin{equation}
\label{eq:ebitda}
\text{EBITDA} = \text{Betriebsergebnis} + \text{Abschreibung} - \text{Wertaufholung}
\end{equation}

Das ergibt \MioUSDexakt{\EBITDA} nach \MioUSDexakt{\EBITDAVJ} und eine Marge
von \Pct{\EBITDAMarge} nach \Pct{\EBITDAMargeVJ}.

\subsection{Konzernergebnis}

Das Konzernergebnis stieg von \MioUSDexakt{\KonzernergebnisVJ} auf
\MioUSDexakt{\Konzernergebnis}, also um \Pct{\ErgebnisDelta}.\quelleFS{2025}{7}\quelleMerken{s8fsxcacfxh}
Diese Zahl ist ohne die beiden Einmaleffekte nicht zu lesen: die
Wertaufholung von \MioUSDexakt{\Wertaufholung} und ein
Ver\"au{\ss}erungsgewinn von \MioUSDexakt{\VeraeusserungsGewinn}.

Das vom Unternehmen selbst berichtete \emph{bereinigte} Ergebnis betr\"agt
\MioUSDexakt{\BereinigtesErgebnis} nach
\MioUSDexakt{\BereinigtesErgebnisVJ}, ein Zuwachs von
\Pct{\BereinigtDelta}.\quelleMDA{2025}{4}\quelleMerken{s8mdaxcacfxe}

\bk{Der Unterschied ist erheblich und f\"ur das Anlageurteil entscheidend:
Das berichtete Ergebnis ist um mehr als die H\"alfte gr\"o{\ss}er als das
bereinigte. Teil~\ref{sec:flow} legt die H\"urde deshalb auf das bereinigte
Ergebnis. Eine einmalige Bilanzkorrektur gibt keinen Ma{\ss}stab f\"ur eine
Daueranlage ab.}

\subsection{Bilanz und Nettoliquidit\"at}
\label{sec:bilanz}

Die Bilanzsumme betr\"agt \MioUSDexakt{\Bilanzsumme} nach
\MioUSDexakt{\BilanzsummeVJ}, das Eigenkapital \MioUSDexakt{\Eigenkapital}
nach \MioUSDexakt{\EigenkapitalVJ}; die Eigenkapitalquote liegt damit bei
\Pct{\Eigenkapitalquote}.\quelleFS{2025}{6}\quelleMerken{s8fsxcacfxg}

Eine Unterscheidung ist hier tragend. Die Gesamtverbindlichkeiten betragen
\MioUSDexakt{\Gesamtverbindlichkeiten}, die \emph{Finanzschulden} dagegen nur
\MioUSDexakt{\Finanzschulden}; der Unterschied besteht \"uberwiegend aus
latenten Steuern und Rekultivierungsverpflichtungen. Den ersten Wert f\"ur
den zweiten zu nehmen verkehrt die Verschuldungslage ins Gegenteil, denn dem
Kassenbestand von \MioUSDexakt{\Liquiditaet} stehen Finanzschulden von
\MioUSDexakt{\Finanzschulden} gegen\"uber:

\begin{equation}
\label{eq:nettoliquiditaet}
\text{Nettoliquidit\"at} = \text{Kasse} - \text{Finanzschulden}
\end{equation}

Sie betr\"agt \MioUSDexakt{\Nettoliquiditaet} nach
\MioUSDexakt{\NettoliquiditaetVJ}. Von der revolvierenden Kreditlinie sind
\MioUSDexakt{\FazilitaetUnbeansprucht} unbeansprucht.\quelleFS{2025}{31}\quelleMerken{s8fsxcacfxdb}

\bk{Alamos hat keine Verschuldung, die eine Entschuldungsrechnung
rechtfertigte. Die Frage, die bei einem verschuldeten Unternehmen zuerst zu
stellen w\"are, stellt sich hier nicht; an ihre Stelle tritt in
Teil~\ref{sec:flow} die Finanzierungsfrage.}

\subsection{Dividende}

Erkl\"art wurden Dividenden von \MioUSDexakt{\DividendeErklaert}, davon
\MioUSDexakt{\DividendeGezahlt} in bar.\quelleFS{2025}{37}\quelleMerken{s8fsxcacfxdh} Am
\Datum{18.\,Februar 2026} wurde die Quartalsdividende um
\Pct{\DividendeErhoehung} auf \USDje{\DividendeQuartalNeu} je Aktie
angehoben.\quelleErneut{s8fsxcacfxdh} Bezogen auf den Kurs vom \KursStand{}
ergibt das eine Dividendenrendite von \Pct{\Dividendenrendite}.

\subsection{Die Zehnjahresreihe und die Lage des j\"ungsten Wertes}
\label{sec:reihe}

Zwei Jahre gen\"ugen, um zu sagen, was sich ge\"andert hat. Sie gen\"ugen
nicht, um zu wissen, ob der j\"ungste Wert normal ist -- und daran h\"angt
das Urteil. Tabelle~\ref{tab:reihe} f\"uhrt deshalb zehn Jahre.

\begin{table}[htbp]
\centering
\caption{Zehnjahresreihe 2016--2025. Jede Zahl stammt aus dem Bericht ihres
eigenen Jahres, nicht aus der Vergleichsspalte eines sp\"ateren. Betr\"age in
Mio.\,USD, Kosten in USD je Unze. Der freie Cash-Flow ist gerechnet als
operativer Cash-Flow abz\"uglich Sachinvestitionen und in allen zehn Jahren
gleich abgegrenzt.}
\label{tab:reihe}
\small
\begin{tabular}{lrrrrrrr}
\toprule
Jahr & Umsatz & Produktion & AISC & Op.\ CF & Investition & Freier CF & Eigenkapital\\
\midrule
\tabellenkoerper{data/reihe_tabelle}
\bottomrule
\end{tabular}
\end{table}

\FloatBarrier

Der Median des freien Cash-Flows der zehn Jahre betr\"agt
\MioUSDexakt{\FCFMedian}. Der Wert des Jahres 2025 betr\"agt
\MioUSDexakt{\FCFReiheXXV} und ist der h\"ochste der Reihe; er liegt im
\Pct{\FCFPerzentil}-Perzentil der eigenen Geschichte. Die Produktion liegt im
\Pct{\ProduktionPerzentil}-Perzentil, die St\"uckkosten im
\Pct{\AISCPerzentil}-Perzentil.

\bk{Das ist der wichtigste Befund dieses Teils, und er schneidet gegen den
Fall. Alamos hat in sieben der zehn Jahre praktisch keinen freien Cash-Flow
erwirtschaftet: Der Median liegt bei \MioUSDexakt{\FCFMedian}, der kleinste
Wert bei \MioUSDexakt{\FCFMin}. Erst die letzten drei Jahre sind deutlich
positiv, und sie fallen mit einer Verdopplung des Goldpreises zusammen.

Eine Bewertung, die auf dem Jahr 2025 aufsetzt, setzt damit auf dem besten
Jahr auf, das dieses Unternehmen je hatte, beim h\"ochsten Goldpreis, den es
je erzielt hat -- und zugleich bei der zweith\"ochsten Kostenbelastung der
Reihe. Wer das nicht ausspricht, liest ein Zyklushoch als Normalfall.
Teil~\ref{sec:flow} rechnet deshalb drei vom Unternehmen selbst berichtete
Goldpreise durch und nicht einen.}

\subsection*{Kennzahlen\"uberblick}
\addcontentsline{toc}{subsection}{Kennzahlen\"uberblick}

\begin{table}[htbp]
\centering
\caption{Kennzahlen\"uberblick 2024 gegen 2025.}
\label{tab:ueberblick}
\begin{tabular}{lrrr}
\toprule
Kennzahl & 2024 & 2025 & Ver\"anderung\\
\midrule
Umsatz (Mio.\,USD) & \MioUSDexakt{\UmsatzVJ} & \MioUSDexakt{\Umsatz} & \Pct{\UmsatzDelta}\\
Produktion (Unzen) & \num{\ProduktionXXIV} & \num{\ProduktionXXV} & \Pct{\ProduktionDelta}\\
Erzielter Goldpreis & \JeUnze{\GoldpreisVJ} & \JeUnze{\Goldpreis} & \Pct{\GoldpreisDelta}\\
AISC & \JeUnze{\AISCVJneu} & \JeUnze{\AISCXXV} & \Pct{\AISCDelta}\\
Marge je Unze & \JeUnze{\MargeJeUnzeVJ} & \JeUnze{\MargeJeUnze} & \Pct{\MargeJeUnzeDelta}\\
EBITDA (Mio.\,USD) & \MioUSDexakt{\EBITDAVJ} & \MioUSDexakt{\EBITDA} & --\\
Konzernergebnis (Mio.\,USD) & \MioUSDexakt{\KonzernergebnisVJ} & \MioUSDexakt{\Konzernergebnis} & \Pct{\ErgebnisDelta}\\
Bereinigtes Ergebnis (Mio.\,USD) & \MioUSDexakt{\BereinigtesErgebnisVJ} & \MioUSDexakt{\BereinigtesErgebnis} & \Pct{\BereinigtDelta}\\
Operativer Cash-Flow (Mio.\,USD) & \MioUSDexakt{\OpCashflowVJ} & \MioUSDexakt{\OpCashflow} & \Pct{\OpCashflowDelta}\\
Freier Cash-Flow (Mio.\,USD) & \MioUSDexakt{\FCFReiheXXIV} & \MioUSDexakt{\FCFReiheXXV} & --\\
Nettoliquidit\"at (Mio.\,USD) & \MioUSDexakt{\NettoliquiditaetVJ} & \MioUSDexakt{\Nettoliquiditaet} & --\\
Eigenkapital (Mio.\,USD) & \MioUSDexakt{\EigenkapitalVJ} & \MioUSDexakt{\Eigenkapital} & --\\
\bottomrule
\end{tabular}
\end{table}
